\chapter{Cách đọc cuốn sách này}

Cuốn sách này dạy mọi công nghệ có trong dự án TeacherSupporter ---
Spring Boot và Spring Cloud, Postgres và MongoDB, Kafka, Docker,
Kubernetes, và pipeline Jenkins gắn kết tất cả lại với nhau --- bằng
chính các file của dự án làm ví dụ. Không có gì được bịa ra cho trang
sách: mọi cấu hình bạn sẽ đọc đều đang được triển khai, hoặc sắp được
triển khai, trên một máy chủ gia đình thật hai nhân tên là
\code{ts-server}.

\section*{Những quy tắc cuốn sách tuân theo}

\begin{description}
  \item[Khái niệm trước, cấu hình sau.] Mỗi chương trước hết dạy ý tưởng
  bằng ngôn ngữ trung lập, không gắn với nhà cung cấp nào, rồi đọc file
  thật của dự án từng dòng một, rồi đưa một hai ví dụ từ \emph{bên ngoài}
  dự án để bạn phân biệt được cái gì là quy luật chung, cái gì chỉ là thói
  quen cục bộ.
  \item[Quyết định là nội dung chính.] Ở bất cứ đâu dự án phải chọn giữa
  các phương án (Jib thay vì \code{docker build}, Kustomize thay vì Helm,
  một server Postgres thứ hai thay vì một cơ sở dữ liệu thứ hai), lý do
  nằm trong hộp \textsc{decision}. Đó chính là những câu hỏi ``tại sao
  bạn lại\ldots'' mà các buổi phỏng vấn được dựng nên từ đó.
  \item[Lỗi được trình bày từ triệu chứng.] Các hộp \textsc{pitfall} bắt
  đầu bằng thứ bạn thực sự nhìn thấy (\code{CrashLoopBackOff}, một lỗi
  TLS, một 404 lặng lẽ) rồi lần ngược về nguyên nhân, vì đó là hướng mà
  việc gỡ lỗi thật diễn ra.
\end{description}

\section*{Chú giải các hộp}

\begin{decision}[ví dụ]
Vì sao dự án chọn thiết kế này thay vì thiết kế kia, và điều gì phải thay
đổi để lựa chọn còn lại thắng thế.
\end{decision}

\begin{pitfall}[ví dụ]
Một triệu chứng sớm muộn bạn cũng sẽ gặp, kèm nguyên nhân và cách sửa.
\end{pitfall}

\begin{handson}[ví dụ]
Một lệnh để chạy trên stack thật, cùng kết quả bạn nên mong đợi. Cuốn
sách này được viết để đọc bên cạnh một terminal.
\end{handson}

\begin{hardware}[ví dụ]
Một lựa chọn chỉ đúng trên chiếc máy này (2 nhân, 19.4\,GB RAM, SSD
SATA, Wi-Fi). Đừng bê những thứ này vào trung tâm dữ liệu.
\end{hardware}

\begin{quiz}
Mỗi chương kết thúc bằng vài câu hỏi trong hộp này. Đáp án ở
Phụ lục~A. Hãy viết câu trả lời ra giấy trước khi xem đáp án --- nhận ra
không phải là nhớ được.
\end{quiz}

\section*{Ký hiệu}

Code inline và tên file trông như thế này: \code{application.yml}. Các
listing có chú thích mang những dấu khoanh tròn --- \co{1}, \co{2} --- được
giải thích ngay bên dưới listing. Lệnh shell được hiển thị với dấu nhắc
\code{\$} cho server và \code{PS>} cho Windows PowerShell; đừng gõ dấu
nhắc.

\section*{Điều gì được giả định, điều gì không}

Giả định: bạn đọc được Java, và đã dùng Maven và git. Không giả định:
bất kỳ kiến thức nào trước đó về Kubernetes, Jenkins, Kafka, bên trong
Docker, hay frontend. Ở những chỗ tài liệu của dự án đã trình bày việc cài
đặt máy một lần (\code{docs/RUNBOOK-SERVER-SETUP.md}), cuốn sách chỉ tóm
tắt rồi đi tiếp.
